pdfoutput=1
\pdfoutput=1
%=============================================================================
% Supplementary Information S1
% Theorem 1: Multi-Expert Consistency Noise Detection Guarantee
% Nature SI — LaTeX
%=============================================================================

\section{Theorem 1: Multi-Expert Consistency Noise Detection Guarantee (The Unlikely Lone Genius Theorem)}
\label{sec:thm1}

We establish a formal guarantee for the SCX noise detection procedure. When $M$
experts, each trained on a disjoint data subset, exhibit high consensus on a
sample's erroneousness, that sample can be identified as label noise with
confidence that converges exponentially fast to $1$ as $M$ grows.

%------------------------------------------------------------------------------
\subsection{Notation and Setup}
\label{sec:thm1:setup}

Let $(\mathcal{X},\mathcal{Y})$ be the input--label space for
$0$-$1$ classification with $|\mathcal{Y}| = K_{\mathcal{Y}}$.
Data follows $(X,Y) \sim \mathcal{D}$ with true labeling function
$f^* : \mathcal{X} \to \mathcal{Y}$ (unknown oracle). We have $M$ expert models
$\{f_m\}_{m=1}^M$, each $f_m : \mathcal{X} \to \mathcal{Y}$.

The state space $\mathcal{S}$ indexes a measurable partition
$\Pi = \{s_1,\dots,s_{|\mathcal{S}|}\}$ of $\mathcal{X}$, where each state
$s \in \mathcal{S}$ groups samples with similar characteristics.

\begin{definition}[Label noise model]
\label{def:noise-model}
For a sample $(x, y^*)$ with true label $y^* = f^*(x)$, the observed label $y$ is
generated as
\begin{equation}
y = \begin{cases}
y^* & \text{with probability } 1 - \eta, \\[4pt]
\operatorname{Uniform}(\mathcal{Y} \setminus \{y^*\}) & \text{with probability } \eta,
\end{cases}
\label{eq:noise-model}
\end{equation}
where $\eta \in (0, \tfrac12)$ is the global noise rate and the noise event is
independent of $x$ and of all training data.
\end{definition}

Let $\ell : \mathcal{Y} \times \mathcal{Y} \to \{0,1\}$ be the $0$-$1$
classification loss, and $\tau > 0$ an error threshold. Define the
expert error indicator and the consensus score as
\begin{align}
e_m(x, y) &= \mathbf{1}\{\ell(f_m(x), y) > \tau\}, &
C(x) &= \frac{1}{M}\sum_{m=1}^M e_m(x, y).
\label{eq:consensus-score}
\end{align}

\begin{definition}[Detection rule]
\label{def:detection-rule}
A sample $(x,y)$ is flagged as noisy if and only if $C(x) > \theta$, where
$\theta \in (0,1)$ is a detection threshold.
\end{definition}

%------------------------------------------------------------------------------
\subsection{Assumptions (A1)--(A6)}
\label{sec:thm1:assumptions}

The following assumptions are used throughout Theorem~1 and its supporting
lemmas.

\begin{enumerate}
\item[\textbf{(A1)}] \textbf{(Disjoint training sets).}
  The $M$ experts are trained on $M$ disjoint independent and identically
  distributed subsets:
  \begin{equation}
  D_m \sim \mathcal{D}^{n_m},\quad D_m \cap D_{m'} = \varnothing,\quad
  D_m \perp D_{m'} \;\text{ for } m \neq m'.
  \label{eq:a1}
  \end{equation}

\item[\textbf{(A2)}] \textbf{(Conditional independence on clean data).}
  For any clean sample $(x,y)$ (i.e.\ $y = y^*$), the error indicators
  $\{e_m(x,y)\}_{m=1}^M$ are conditionally independent given $x$.
  \begin{remark}
  This follows from (A1): $e_m(x,y) = \mathbf{1}\{\ell(f_m(x),y) > \tau\}$
  depends only on $f_m$, and $f_m$ is a function of $D_m$; since the $D_m$
  are mutually independent, $\{f_m(x)\}_m$ and consequently
  $\{e_m(x,y)\}_m$ are conditionally independent given $x$.
  \end{remark}

\item[\textbf{(A3)}] \textbf{(0--1 classification loss).}
  $\ell(a,b) \in \{0,1\}$ is the $0$-$1$ loss for all $a,b \in \mathcal{Y}$.

\item[\textbf{(A4)}] \textbf{(Uniform independent noise).}
  The label-flipping event is independent of $x$ and of all $D_m$. The noise
  label is uniformly distributed over $\mathcal{Y} \setminus \{y^*\}$.

\item[\textbf{(A5)}] \textbf{(State homogeneity).}
  The state partition $\Pi$ is such that within each state $s$, the clean-data
  error rate of the experts is approximately uniform. Formally, there exist
  state-level constants $\{\mu_s\}_{s \in \mathcal{S}}$ satisfying
  \begin{equation}
  \sup_{x \in s} \; \mathbb{E}[C(X) \mid \text{clean}, X = x] \leq \mu_s,
  \qquad \forall s \in \mathcal{S}.
  \label{eq:a5}
  \end{equation}

\item[\textbf{(A6)}] \textbf{(Balanced error distribution).}
  The experts' error probabilities do not concentrate excessively on any
  single incorrect class. There exists a constant $C_{\text{bal}} \geq 1$
  such that for any state $s$ and any $x \in s$,
  \begin{equation}
  \max_{c \neq y^*} \mu_c(x) \;\leq\; C_{\text{bal}} \cdot
  \frac{\mu_s}{K_{\mathcal{Y}} - 1},
  \qquad
  \mu_c(x) = \frac{1}{M}\sum_{m=1}^M \mathbb{P}(f_m(x) = c \mid x).
  \label{eq:a6}
  \end{equation}
  The case $C_{\text{bal}} = 1$ corresponds to perfectly uniform errors
  (every incorrect class equally likely). $C_{\text{bal}} = 2$ is a
  conservative default. This assumption is empirically testable.
\end{enumerate}

%------------------------------------------------------------------------------
\subsection{Lemma 1: Mean Separation}
\label{sec:thm1:lemma1}

\begin{lemma}[Mean Separation]
\label{lem:mean-separation}
Under assumptions (A1)--(A5), for any $x \in s$,
\begin{align}
\mathbb{E}[C(X) \mid \text{clean}, X = x] &\leq \mu_s,
\label{eq:lem1-clean} \\
\mathbb{E}[C(X) \mid \text{noise}, X = x]
&= 1 - \frac{1}{K_{\mathcal{Y}} - 1} \,
   \mathbb{E}[C(X) \mid \text{clean}, X = x]
   \;\geq\; 1 - \frac{\mu_s}{K_{\mathcal{Y}} - 1}.
\label{eq:lem1-noise}
\end{align}

Consequently, when $\mu_s < \dfrac{K_{\mathcal{Y}} - 1}{K_{\mathcal{Y}}}$
there exists $\theta$ such that
\[
\mathbb{E}[C \mid \text{clean}] < \theta < \mathbb{E}[C \mid \text{noise}],
\]
and the separation gap for state $s$ is
\begin{equation}
\Delta_s(\theta)
= \min\!\Bigl(\theta - \mu_s,\;
  1 - \frac{\mu_s}{K_{\mathcal{Y}}-1} - \theta\Bigr).
\label{eq:separation-gap}
\end{equation}

With the optimal threshold
\[
\theta_s^* = \frac12\!\left(1 + \mu_s\cdot\frac{K_{\mathcal{Y}}-2}{K_{\mathcal{Y}}-1}\right)\!,
\]
the maximised gap (in the ideal case $C_{\text{bal}}=1$) is
\[
\Delta_s^* = \frac12\!\left(1 - \mu_s\cdot\frac{K_{\mathcal{Y}}}{K_{\mathcal{Y}}-1}\right).
\]
When $C_{\text{bal}} > 1$, the optimal threshold shifts towards $\mu_s$ to
compensate for the non-uniform error distribution.
\end{lemma}

\begin{proof}[Proof of Lemma~1]
The clean-sample bound follows directly from (A5). For a noise sample,
\begin{align*}
\mathbb{E}[C \mid \text{noise}, X = x]
&= \frac{1}{M}\sum_{m=1}^M \mathbb{P}\bigl(\ell(f_m(x), y) > \tau
   \;\big|\; \text{noise}, x\bigr) \\
&= \frac{1}{M}\sum_{m=1}^M \mathbb{P}\bigl(f_m(x) \neq y
   \;\big|\; \text{noise}, x\bigr) \qquad(\text{for } 0\text{-}1\text{ loss}) \\
&= 1 - \frac{1}{M}\sum_{m=1}^M \mathbb{P}\bigl(f_m(x) = y
   \;\big|\; \text{noise}, x\bigr).
\end{align*}
By (A4), the noise label $y$ is uniform over $\mathcal{Y}\setminus\{y^*\}$ and
independent of all experts:
\begin{align*}
\mathbb{P}(f_m(x) = y \mid \text{noise}, x)
&= \sum_{c \neq y^*} \mathbb{P}(y = c \mid \text{noise})\,
   \mathbb{P}(f_m(x) = c \mid x) \\
&= \frac{1}{K_{\mathcal{Y}}-1} \sum_{c \neq y^*}
   \mathbb{P}(f_m(x) = c \mid x) \\
&= \frac{1}{K_{\mathcal{Y}}-1}\,
   \mathbb{P}(f_m(x) \neq y^* \mid x)
 = \frac{1}{K_{\mathcal{Y}}-1}\,
   \mathbb{E}[e_m \mid \text{clean}, x].
\end{align*}
Averaging over $m$ gives
\begin{align*}
\mathbb{E}[C \mid \text{noise}, x]
&= 1 - \frac{1}{K_{\mathcal{Y}}-1}\,
   \frac{1}{M}\sum_{m=1}^M \mathbb{E}[e_m \mid \text{clean}, x] \\
&= 1 - \frac{1}{K_{\mathcal{Y}}-1}\,
   \mathbb{E}[C \mid \text{clean}, x]
\;\geq\; 1 - \frac{\mu_s}{K_{\mathcal{Y}}-1},
\end{align*}
where the inequality uses (A5). The separation gap exists whenever
$\mu_s < (K_{\mathcal{Y}}-1)/K_{\mathcal{Y}}$, because then
$1 - \mu_s/(K_{\mathcal{Y}}-1) > \mu_s$. \end{proof}

%------------------------------------------------------------------------------
\subsection{Lemma 2: False Positive Rate Bound}
\label{sec:thm1:lemma2}

\begin{lemma}[False Positive Rate Upper Bound]
\label{lem:fpr-bound}
For any state $s \in \mathcal{S}$ satisfying $\mu_s < \theta$, the probability
that a clean sample is incorrectly flagged as noise obeys
\begin{equation}
\mathbb{P}\bigl(C > \theta \mid \text{clean}, X \in s\bigr)
\;\leq\; \exp\!\bigl(-2M(\theta - \mu_s)^2\bigr).
\label{eq:fpr-bound}
\end{equation}
\end{lemma}

\begin{proof}
Fix $x \in s$. By (A5), $\mathbb{E}[C \mid \text{clean}, X = x] \leq \mu_s$.
By (A1)--(A2), $\{e_m\}$ are conditionally independent given $x$, and by
(A3) each $e_m \in [0,1]$. Apply Hoeffding's inequality
\cite{hoeffding1963} to $C = \frac{1}{M}\sum_{m} e_m$:
\[
\mathbb{P}\bigl(C - \mathbb{E}[C] > \theta - \mu_s \mid \text{clean}, x\bigr)
\;\leq\; \exp\!\bigl(-2M(\theta - \mu_s)^2\bigr),
\]
where $\theta - \mu_s > 0$ by hypothesis. The bound holds for every $x \in s$,
so integrating over $x$ yields
\begin{align*}
\mathbb{P}(C > \theta \mid \text{clean}, X \in s)
&= \mathbb{E}_{X \mid s}\bigl[ \mathbb{P}(C > \theta \mid \text{clean}, X) \bigr] \\
&\leq \sup_{x \in s} \mathbb{P}(C > \theta \mid \text{clean}, x)
\;\leq\; \exp\!\bigl(-2M(\theta - \mu_s)^2\bigr). \qedhere
\end{align*}
\end{proof}

%------------------------------------------------------------------------------
\subsection{Lemma 3: True Positive Rate Bound}
\label{sec:thm1:lemma3}

\begin{lemma}[True Positive Rate Lower Bound]
\label{lem:tpr-bound}
Under assumptions (A1)--(A6), for any state $s \in \mathcal{S}$ satisfying
\[
\theta < 1 - C_{\text{bal}}\cdot\frac{\mu_s}{K_{\mathcal{Y}}-1},
\]
the probability that a noisy sample is correctly detected obeys
\begin{equation}
\mathbb{P}\bigl(C > \theta \mid \text{noise}, X \in s\bigr)
\;\geq\; 1 - \exp\!\Bigl(
  -2M\Bigl(1 - C_{\text{bal}}\cdot\frac{\mu_s}{K_{\mathcal{Y}}-1}
  - \theta\Bigr)^{\!2}\Bigr).
\label{eq:tpr-bound}
\end{equation}
\end{lemma}

\begin{proof}
Fix $x \in s$ and condition on a specific noise label $c \neq y^*$. For
$0$-$1$ loss, expert $m$ errs when $f_m(x) \neq c$, so
$e_m = \mathbf{1}\{f_m(x) \neq c\}$. By (A1) the $\{f_m\}$ are independent
functions of disjoint data; by (A2) the $\{e_m\}$ are conditionally independent
given $(x,c)$. The conditional expectation is
\[
\mathbb{E}[e_m \mid x, c] = 1 - \mathbb{P}(f_m(x) = c \mid x)
= 1 - \mu_{c,m}(x),
\]
hence
\[
\mathbb{E}[C \mid x, c]
= \frac{1}{M}\sum_{m=1}^M \mathbb{E}[e_m \mid x, c]
= 1 - \mu_c(x)
\;\geq\; 1 - C_{\text{bal}}\cdot\frac{\mu_s}{K_{\mathcal{Y}}-1},
\]
where the inequality follows from (A6). Since $e_m \in [0,1]$ are conditionally
independent, Hoeffding's inequality gives
\begin{align*}
\mathbb{P}(C \leq \theta \mid x, c)
&= \mathbb{P}\!\left(\frac{1}{M}\sum_{m=1}^M e_m \leq \theta
   \;\Big|\; x, c\right) \\
&\leq \exp\!\Bigl(-2M\bigl(1 - \mu_c(x) - \theta\bigr)^2\Bigr) \\
&\leq \exp\!\Bigl(-2M\bigl(1 - C_{\text{bal}}\cdot\frac{\mu_s}{K_{\mathcal{Y}}-1}
   - \theta\bigr)^2\Bigr).
\end{align*}
The Hoeffding gap condition $1 - \mu_c(x) > \theta$ holds because
\[
1 - \mu_c(x) \geq 1 - C_{\text{bal}}\cdot\frac{\mu_s}{K_{\mathcal{Y}}-1}
> \theta
\]
by the lemma's hypothesis and (A6). Averaging over the uniform noise label
$c \in \mathcal{Y}\setminus\{y^*\}$,
\[
\mathbb{P}(C \leq \theta \mid \text{noise}, x)
= \frac{1}{K_{\mathcal{Y}}-1}\sum_{c \neq y^*}
  \mathbb{P}(C \leq \theta \mid x, c)
\;\leq\; \exp\!\Bigl(-2M\bigl(1 - C_{\text{bal}}\cdot\frac{\mu_s}{K_{\mathcal{Y}}-1}
  - \theta\bigr)^2\Bigr).
\]
Therefore $\mathbb{P}(C > \theta \mid \text{noise}, x)
\geq 1 - \exp(-2M(1 - C_{\text{bal}}\cdot\mu_s/(K_{\mathcal{Y}}-1)
- \theta)^2)$. Taking expectations over $X \in s$ completes the proof.
\end{proof}

\begin{remark}
When $C_{\text{bal}} = 1$ (perfectly uniform error distribution), Lemma~3
recovers the optimal bound. The constant $C_{\text{bal}}$ can be estimated from
data and used to calibrate the detection threshold $\theta$.
\end{remark}

%------------------------------------------------------------------------------
\subsection{Proof of Theorem 1: F1 Lower Bound}
\label{sec:thm1:proof}

\begin{theorem}[SCX Noise Detection Guarantee]
\label{thm:noise-detection}
Let assumptions (A1)--(A6) hold. Let $\rho_s = \mathbb{P}(X \in s)$ be the
state probability. For any threshold $\theta$ satisfying, for a given state $s$,
\[
\mu_s < \theta < 1 - C_{\text{bal}}\cdot\frac{\mu_s}{K_{\mathcal{Y}}-1},
\]
define the state-level separation gap
\begin{equation}
\Delta_s \;=\;
\min\!\Bigl(\theta - \mu_s,\;
  1 - C_{\text{bal}}\cdot\frac{\mu_s}{K_{\mathcal{Y}}-1} - \theta\Bigr)
> 0.
\label{eq:separation-gap-cbal}
\end{equation}
When $C_{\text{bal}} = 1$, $\Delta_s$ reduces to the definition in
Lemma~\ref{lem:mean-separation}; when $C_{\text{bal}} > 1$ the noise-side gap
narrows because errors may not be perfectly uniform.

Then the SCX noise detector achieves the F1 lower bound
\begin{equation}
\boxed{\;
\operatorname{F1} \;\geq\; 1 - \frac{1}{\eta}
\sum_{s \in \mathcal{S}} \rho_s \cdot
\exp\!\bigl(-2M\Delta_s^2\bigr)
\;}
\label{eq:f1-bound-hoeffding}
\end{equation}
or equivalently
\begin{equation}
\operatorname{F1} \;\geq\; 1 - \sum_{s \in \mathcal{S}} \rho_s \Bigl[
\exp\!\bigl(-2M\Delta_s^2\bigr)
+ \frac{1-\eta}{\eta}\exp\!\bigl(-2M\Delta_s^2\bigr)
\Bigr].
\label{eq:f1-bound-explicit}
\end{equation}
\end{theorem}

\begin{remark}[Effective number of experts under correlation]
\label{rem:Meff}
The bound above uses the nominal number of experts $M$, which assumes
perfect conditional independence (A2'). When experts share the same
architecture or training paradigm, their errors may exhibit positive
pairwise correlation $\bar{\rho}$. In this case the effective number
of independent experts is reduced to the
\textbf{The Expert Dilution Formula (Expert Dilution Formula)}:
$M_{\text{eff}} = M / (1 + (M-1)\bar{\rho})$
(see Liang \& Zeger, 1986, for the generalized estimating equations
variance-inflation factor). The bound then reads
$\operatorname{F1} \geq 1 - \frac{1}{\eta}\sum_s \rho_s
\exp(-2M_{\text{eff}}\Delta_s^2)$, which remains exponentially
convergent in $M$ but at a rate slowed by the factor $(1+(M-1)\bar{\rho})^{-1}$.
Consistent estimation of $\bar{\rho}$ via jackknife or bootstrap is
standard; see \S\ref{sec:correlation-estimation} for details.
\end{remark}

\begin{proof}
Let $R(x) = \mathbf{1}\{C(x) > \theta\}$ be the detection rule. Define
\begin{align*}
\operatorname{TPR}_s &=
\mathbb{P}(R = 1 \mid \text{noise}, X \in s), \\
\operatorname{FPR}_s &=
\mathbb{P}(R = 1 \mid \text{clean}, X \in s).
\end{align*}
By Lemma~\ref{lem:tpr-bound} and Lemma~\ref{lem:fpr-bound},
\begin{align*}
\operatorname{TPR}_s &\geq 1 -
\exp\!\Bigl(-2M\bigl(1 - C_{\text{bal}}\cdot\frac{\mu_s}{K_{\mathcal{Y}}-1}
- \theta\bigr)^2\Bigr), \\
\operatorname{FPR}_s &\leq
\exp\!\bigl(-2M(\theta - \mu_s)^2\bigr).
\end{align*}

From (A4), the noise event is independent of $X$, so
$\mathbb{P}(\text{noise} \mid X \in s) = \eta$ for all $s$. Aggregating across
states,
\[
\operatorname{TPR} = \sum_s \rho_s \operatorname{TPR}_s,\qquad
\operatorname{FPR} = \sum_s \rho_s \operatorname{FPR}_s.
\]

The F1 score is
\[
\operatorname{F1}
= \frac{2\eta\operatorname{TPR}}{2\eta\operatorname{TPR}
   + (1-\eta)\operatorname{FPR} + \eta(1-\operatorname{TPR})}
= \frac{2\eta\operatorname{TPR}}{\eta(1+\operatorname{TPR})
   + (1-\eta)\operatorname{FPR}}.
\]

Set
\begin{align*}
\delta_1 &= \sum_{s} \rho_s \,
\exp\!\Bigl(-2M\bigl(1 - C_{\text{bal}}\cdot\frac{\mu_s}{K_{\mathcal{Y}}-1}
- \theta\bigr)^2\Bigr), \\
\delta_2 &= \sum_{s} \rho_s \,
\exp\!\bigl(-2M(\theta - \mu_s)^2\bigr).
\end{align*}
Then $\operatorname{TPR} \geq 1 - \delta_1$ and
$\operatorname{FPR} \leq \delta_2$. Substituting,
\begin{align*}
\operatorname{F1}
&\geq \frac{2\eta(1-\delta_1)}{\eta(2-\delta_1) + (1-\eta)\delta_2} \\
&= 1 - \frac{\eta\delta_1 + (1-\eta)\delta_2}
         {\eta(2-\delta_1) + (1-\eta)\delta_2} \\
&\geq 1 - \frac{\eta\delta_1 + (1-\eta)\delta_2}{\eta}
\qquad(\text{since denominator }\geq \eta) \\
&= 1 - \delta_1 - \frac{1-\eta}{\eta}\,\delta_2 .
\end{align*}

Now observe that $\Delta_s$ as defined in
\eqref{eq:separation-gap-cbal} satisfies
\[
\exp\!\bigl(-2M(\theta - \mu_s)^2\bigr) \leq \exp\!\bigl(-2M\Delta_s^2\bigr),
\quad
\exp\!\Bigl(-2M\bigl(1 - C_{\text{bal}}\cdot\frac{\mu_s}{K_{\mathcal{Y}}-1}
- \theta\bigr)^2\Bigr) \leq \exp\!\bigl(-2M\Delta_s^2\bigr).
\]
Therefore
\begin{align*}
\operatorname{F1}
&\geq 1 - \sum_s \rho_s\Bigl[
   \exp\!\bigl(-2M\Delta_s^2\bigr)
   + \frac{1-\eta}{\eta}\exp\!\bigl(-2M\Delta_s^2\bigr)
\Bigr] \\
&= 1 - \frac{1}{\eta}\sum_s \rho_s\exp\!\bigl(-2M\Delta_s^2\bigr).
\qedhere
\end{align*}
\end{proof}

\begin{corollary}[Asymptotic optimality]
\label{cor:asymptotic}
As $M \to \infty$, for all states satisfying
$\mu_s < (K_{\mathcal{Y}}-1)/K_{\mathcal{Y}}$,
\[
\operatorname{F1} = 1 - \mathcal{O}_P\!\left(
\frac{1}{\eta}\, e^{-2M\Delta_{\min}^2}\right),
\qquad
\Delta_{\min} = \min_{s \in \mathcal{S}} \Delta_s.
\]
\end{corollary}

\begin{example}[Numerical illustration]
\label{ex:thm1-numerical}
Take $K_{\mathcal{Y}} = 10$ classes, $\eta = 0.1$ noise rate, $M = 20$ experts,
and per-state clean error bound $\mu_s = 0.2$. With $C_{\text{bal}} = 1$, the
optimal separation gap is
\[
\Delta_s^* = \frac12\!\left(1 - 0.2 \cdot \frac{10}{9}\right)
\approx 0.389,
\]
and the F1 lower bound from Theorem~\ref{thm:noise-detection} is
\[
\operatorname{F1} \;\geq\; 1 - 10 \cdot
\exp\!\bigl(-2 \cdot 20 \cdot 0.389^2\bigr)
= 1 - 10 \cdot e^{-6.05}
> 0.976.
\]
In practice, if the clean expert error rate is higher (e.g.\
$\mu_s \approx 0.45$ as in CIFAR-10 with lightly trained experts), the gap
shrinks to $\Delta_s \approx 0.25$, giving the weaker bound
$\operatorname{F1} \geq 0.18$, which remains valid as a guarantee but is not
tight. This sensitivity to $\mu_s$ underscores the importance of estimating
$\mu_s$ accurately from validation data.
\end{example}

%------------------------------------------------------------------------------
\subsection{Chernoff Tightening}
\label{sec:thm1:chernoff}

Lemmas~\ref{lem:fpr-bound} and \ref{lem:tpr-bound} can be sharpened using the
Chernoff bound instead of Hoeffding. For clean data, let
$\{e_m\}$ be independent Bernoulli variables with mean
$\mu = \mathbb{E}[C \mid \text{clean}, x] \leq \mu_s$. For
$\theta > \mu$,
\begin{align*}
\mathbb{P}(C > \theta \mid \text{clean}, x)
&= \mathbb{P}\!\left(\frac{1}{M}\sum_{m} e_m > \theta\right) \\
&\leq \inf_{\lambda > 0}\, e^{-\lambda M\theta}\,
   \mathbb{E}\!\left[e^{\lambda\sum_m e_m}\right] \\
&= \inf_{\lambda > 0}\, e^{-\lambda M\theta}\,
   \prod_{m=1}^M (1 - \mu_m + \mu_m e^\lambda) \\
&\leq \exp\!\bigl(-M \cdot \operatorname{KL}(\theta \,\|\, \mu_s)\bigr),
\end{align*}
where $\operatorname{KL}(p \,\|\, q) = p\log\frac{p}{q}
+ (1-p)\log\frac{1-p}{1-q}$.

For the noise lower tail, conditioning on label $c$ and applying the same
bound yields
\[
\mathbb{P}(C \leq \theta \mid \text{noise}, x, c)
\;\leq\; \exp\!\Bigl(
  -M \cdot \operatorname{KL}\!\Bigl(\theta \;\Big\|\;
   1 - C_{\text{bal}}\cdot\frac{\mu_s}{K_{\mathcal{Y}}-1}\Bigr)
\Bigr).
\]

\begin{remark}[KL direction]
\label{remark:kl-direction}
The standard Chernoff bound for the lower tail is
$\mathbb{P}(\bar{X} \leq \theta) \leq \exp(-n \cdot \operatorname{KL}(\theta \| p))$
when $\theta < p$. Here the true mean is
$p = 1 - C_{\text{bal}}\cdot\mu_s/(K_{\mathcal{Y}}-1)$ and the threshold
$\theta$ satisfies $\theta < p$ by hypothesis, so the correct KL argument is
$\theta$ (first) and $p$ (second).
\end{remark}

Combining the two tails gives the Chernoff-form F1 bound
\begin{multline}
\operatorname{F1} \;\geq\; 1
- \sum_{s \in \mathcal{S}} \rho_s \,
\exp\!\Bigl(-M \cdot \operatorname{KL}\!\Bigl(\theta \;\Big\|\;
  1 - C_{\text{bal}}\cdot\frac{\mu_s}{K_{\mathcal{Y}}-1}\Bigr)\Bigr) \\
- \frac{1-\eta}{\eta} \sum_{s \in \mathcal{S}} \rho_s \,
\exp\!\bigl(-M \cdot \operatorname{KL}(\theta \,\|\, \mu_s)\bigr).
\label{eq:f1-bound-chernoff}
\end{multline}

The Chernoff bound is tighter than the Hoeffding bound by a factor that grows
with the separation gap. Specifically,
\[
\exp(-2M\Delta^2) \geq \exp(-M \cdot \operatorname{KL}(\mu + \Delta \,\|\, \mu)),
\]
with equality only as $\Delta \to 0$. In the practical regime
$\Delta \approx 0.1$--$0.4$, the Chernoff bound is $2$--$5$ times tighter.

%------------------------------------------------------------------------------
\subsection{Corollaries}
\label{sec:thm1:corollaries}

\subsubsection*{Corollary 1: Symmetric Experts}
If all experts share the same clean-data error rate $\varepsilon$ (symmetric
experts), then $\mu_s = \varepsilon$ for all states,
\[
\Delta = \frac12\!\left(1 - \varepsilon\cdot\frac{K_{\mathcal{Y}}}{K_{\mathcal{Y}}-1}\right),
\]
and the F1 bound simplifies to
\begin{equation}
\operatorname{F1} \;\geq\; 1 - \frac{1}{\eta}\,
\exp\!\left(-\frac{M}{2}\Bigl(1 - \frac{\varepsilon K_{\mathcal{Y}}}{K_{\mathcal{Y}}-1}\Bigr)^2\right).
\label{eq:cor-symmetric}
\end{equation}
For binary classification ($K_{\mathcal{Y}}=2$) this further reduces to
\[
\operatorname{F1} \;\geq\; 1 - \frac{1}{\eta}\,
\exp\!\bigl(-2M(\tfrac12 - \varepsilon)^2\bigr).
\]

\subsubsection*{Corollary 2: Optimal Threshold Selection}
Given noise rate $\eta$, number of classes $K_{\mathcal{Y}}$, number of experts
$M$, and maximum clean error rate $\mu_{\max} = \max_s \mu_s$, the optimal
threshold is
\begin{equation}
\theta^* = \arg\max_{\theta} \min_s \Delta_s(\theta)
= \frac12\!\left(1 + \mu_{\max}\cdot\frac{K_{\mathcal{Y}}-2}{K_{\mathcal{Y}}-1}\right).
\label{eq:optimal-threshold}
\end{equation}
The resulting minimum separation gap is
\[
\Delta_{\min}^* = \frac12\!\left(1 - \mu_{\max}\cdot\frac{K_{\mathcal{Y}}}{K_{\mathcal{Y}}-1}\right),
\]
and the minimum number of experts required to achieve
$\operatorname{F1} \geq 1 - \varepsilon_0$ is
\begin{equation}
M \;\geq\; \frac{1}{2\Delta_{\min}^{*2}}\,
\log\!\left(\frac{1}{\eta\varepsilon_0}\right).
\label{eq:min-experts}
\end{equation}
When controlling error across multiple states simultaneously, replace the
numerator $1$ in the logarithm with $|\mathcal{S}|$.

\subsubsection*{Corollary 3: Uniform Detectability}
If there exists $\delta > 0$ such that
$\mu_s \leq (K_{\mathcal{Y}}-1)/K_{\mathcal{Y}} - \delta$ for all $s$, then
with the fixed threshold $\theta = \tfrac12$,
\[
\Delta_s \geq \frac{\delta}{2},\qquad
\operatorname{F1} \;\geq\; 1 - \frac{1}{\eta}\,
\exp\!\left(-\frac{M\delta^2}{2}\right).
\]
Thus whenever every state's clean expert error rate is uniformly bounded
below $(K_{\mathcal{Y}}-1)/K_{\mathcal{Y}}$, a fixed threshold yields
exponential F1 guarantees.

\subsubsection*{Corollary 4: Finite-Sample Correction}
In practice $\mu_s$ must be estimated from finite validation data. Let
$n_s$ be the number of clean validation samples in state $s$, and let
$\hat{\mu}_s = \frac{1}{M}\sum_m \hat{\varepsilon}_m(s)$ be the empirical
estimate, where $\hat{\varepsilon}_m(s)$ is expert $m$'s empirical error
rate on state $s$. By Hoeffding and the union bound, with probability at
least $1 - \delta_0$,
\[
|\hat{\mu}_s - \mu_s| \;\leq\; \sqrt{\frac{\log(2M/\delta_0)}{2n_s}},
\qquad \forall m, s.
\]
Define the conservative upper bound
\tilde{\mu}_s = \hat{\mu}_s + \sqrt{\log(2M/\delta_0)/(2n_s)}.
Then Theorem~\ref{thm:noise-detection} holds with probability
$\geq 1 - \delta_0$ when $\tilde{\mu}_s$ replaces $\mu_s$ in
$\Delta_s$.

%------------------------------------------------------------------------------
\subsection{Connection to Dawid--Skene}
\label{sec:thm1:dawid-skene}

\subsubsection*{Dawid--Skene model}
Dawid and Skene \cite{dawid1979} estimate true labels and annotator reliability
from multiple independent annotations via EM. Each annotator $m$ is assigned a
global confusion matrix
\[
\pi_m(c' \mid c) = \mathbb{P}(f_m(x) = c' \mid y^* = c), \qquad \forall x \in \mathcal{X}.
\]
The final label estimate is a weighted vote
$\hat{y}_{\text{DS}} = \arg\max_c \sum_m w_m \mathbf{1}\{f_m(x)=c\}$,
where weights $w_m$ are derived from the diagonal entries of $\pi_m$.

\subsubsection*{SCX generalization}
SCX replaces the global confusion matrix with a state-conditioned variant:
\[
\pi_m(c' \mid c, s) = \mathbb{P}(f_m(x) = c' \mid y^* = c,\; x \in s).
\]
SCX weights therefore depend on the state:
\[
w_m(x) = \frac{\exp(-\alpha \hat{R}_m(s(x)))}
              {\sum_{m'} \exp(-\alpha \hat{R}_{m'}(s(x)))},
\qquad \hat{R}_m(s) = \sum_c \pi_m(c \mid c, s).
\]

\subsubsection*{Comparison}
The SCX detection rule in Theorem~1 is built on the \emph{consensus} concept
-- noisy samples exhibit high error across \emph{all} experts -- which is
qualitatively different from Dawid--Skene's global reliability estimation.
Key differences are summarised in Table~\ref{tab:ds-vs-scx}.

\begin{table}[h]
\centering
\caption{Comparison of Dawid--Skene and SCX (Theorem~1) for noise detection.}
\label{tab:ds-vs-scx}
\begin{tabular}{lp{5.5cm}p{5.5cm}}
\toprule
Dimension & Dawid--Skene & SCX (Theorem~1) \\
\midrule
Reliability & Global constant $\pi_m$ & State-conditioned $\pi_m(s)$ \\
Noise detection & Posterior agreement & Multi-expert error consensus \\
Annotations needed & Multiple ($\geq 3$) annotators & Multiple expert models \\
Class imbalance robustness & Weak & Strong \\
Separates noise from hardness & No & Yes (consensus score $C$) \\
Theoretical guarantee & Asymptotic consistency & Exponential convergence \\
\bottomrule
\end{tabular}
\end{table}

When the state partition degenerates to the trivial partition
$\mathcal{S} = \{\mathcal{X}\}$, SCX recovers Dawid--Skene as a special case:
weights become global constants and the consensus score becomes the global
average error rate.

%------------------------------------------------------------------------------
% Local bibliography entries for this section
\begin{thebibliography}{10}
\bibitem{hoeffding1963}
W.~Hoeffding, ``Probability inequalities for sums of bounded random variables,''
\emph{J. Amer. Statist. Assoc.}, vol.~58, no.~301, pp.~13--30, 1963.

\bibitem{dawid1979}
A.~P.~Dawid and A.~M.~Skene, ``Maximum likelihood estimation of observer
error-rates using the EM algorithm,''
\emph{J. R. Statist. Soc. Ser. C}, vol.~28, no.~1, pp.~20--28, 1979.
\end{thebibliography}

\endinput
